% FINIX — Phase 3 (FORTIFY) competition submission.
% Compile: pdflatex -interaction=nonstopmode -halt-on-error docs/PHASE3-SUBMISSION.tex  (run twice for the ToC)
% Every repository path cited here was verified to exist in the repository.
\documentclass[11pt,a4paper]{article}

\usepackage[T1]{fontenc}
\usepackage[utf8]{inputenc}
\usepackage[a4paper,margin=22mm,top=26mm,bottom=24mm]{geometry}
\usepackage{graphicx}
\usepackage{booktabs}
\usepackage{tabularx}
\usepackage{longtable}
\usepackage{array}
\usepackage{xcolor}
\usepackage{enumitem}
\usepackage{fancyhdr}
\usepackage{titlesec}
\usepackage{amsmath}
\usepackage[hidelinks]{hyperref}

% ---------------------------------------------------------------- palette
\definecolor{finixnavy}{RGB}{16,42,74}
\definecolor{finixteal}{RGB}{0,110,120}
\definecolor{finixgrey}{RGB}{96,104,116}
\definecolor{finixrule}{RGB}{205,212,222}
\definecolor{finixbg}{RGB}{243,246,250}

% ---------------------------------------------------------------- headings
\titleformat{\section}{\color{finixnavy}\Large\bfseries}{\thesection}{0.7em}{}
\titleformat{\subsection}{\color{finixnavy}\large\bfseries}{\thesubsection}{0.6em}{}
\titleformat{\subsubsection}{\color{finixteal}\normalsize\bfseries}{}{0em}{}
\titlespacing*{\section}{0pt}{14pt}{6pt}
\titlespacing*{\subsection}{0pt}{11pt}{4pt}

\pagestyle{fancy}
\fancyhf{}
\renewcommand{\headrulewidth}{0.4pt}
\renewcommand{\footrulewidth}{0pt}
\fancyhead[L]{\small\color{finixnavy}\textbf{FINIX} \textcolor{finixgrey}{--- Phase 3 FORTIFY}}
\fancyhead[R]{\small\color{finixgrey}Production Deployment \& Operational Evidence}
\fancyfoot[C]{\small\color{finixgrey}\thepage}

\setlist[itemize]{leftmargin=1.35em,itemsep=1.5pt,topsep=3pt,parsep=0pt}
\setlist[enumerate]{leftmargin=1.6em,itemsep=1.5pt,topsep=3pt,parsep=0pt}
\renewcommand{\arraystretch}{1.25}

% ---------------------------------------------------------------- helpers
\newcommand{\pth}[1]{\texttt{\small #1}}
\newcommand{\chk}{\raisebox{0.4pt}{\framebox[2.6mm]{\rule{0pt}{1.6mm}}}}
\newcommand{\rubric}[1]{{\small\textcolor{finixteal}{\textbf{Rubric:}~#1}}}

% Screenshot placeholder — compiles with no image file present.
\newcommand{\shot}[3]{%
  \begin{center}
  \fbox{%
    \parbox{0.92\textwidth}{%
      \vspace{3pt}\centering
      {\footnotesize\textbf{\textcolor{finixnavy}{SCREENSHOT PLACEHOLDER --- #1}}}\\[3pt]
      {\footnotesize #2}\\[3pt]
      {\footnotesize\itshape\textcolor{finixgrey}{Evidence to capture: #3}}
      \vspace{4pt}}}
  \end{center}}

% Numbered evidence figure placeholder for the evidence appendix.
\newcommand{\evidence}[3]{%
  \begin{center}
  \fbox{%
    \parbox{0.92\textwidth}{%
      \vspace{4pt}\centering
      {\small\textbf{\textcolor{finixnavy}{Figure #1 --- #2}}}\\[4pt]
      \rule{0pt}{26mm}\\
      {\footnotesize\itshape\textcolor{finixgrey}{Evidence to capture: #3}}
      \vspace{5pt}}}
  \end{center}
  \vspace{2pt}}

\begin{document}

% ================================================================ title
\begin{titlepage}
\centering
\vspace*{22mm}
{\color{finixteal}\rule{0.55\textwidth}{1.2pt}}\\[10mm]
{\Huge\bfseries\color{finixnavy} FINIX}\\[5mm]
{\LARGE\color{finixnavy} Resilient Digital Banking Platform}\\[9mm]
{\Large\bfseries\color{finixteal} Phase 3 --- FORTIFY}\\[4mm]
{\large\color{finixnavy} Production Deployment \& Operational Evidence}\\[8mm]
{\large\itshape\color{finixgrey} Deploy It. Automate It. Keep It Running.}\\[10mm]
{\color{finixteal}\rule{0.55\textwidth}{1.2pt}}\\[14mm]

\begin{tabular}{@{}r@{\hspace{1.5em}}l@{}}
\textbf{\color{finixnavy}Team} & [TEAM NAME]\\
\textbf{\color{finixnavy}Competition} & Duothan 6.0\\
\textbf{\color{finixnavy}Date} & [SUBMISSION DATE]\\
\end{tabular}

\vspace{12mm}
\begin{tabular}{@{}r@{\hspace{1.5em}}l@{}}
\textbf{\color{finixnavy}Customer platform} & \url{https://roboti.qzz.io}\\
\textbf{\color{finixnavy}Operations console} & \url{https://admin.roboti.qzz.io}\\
\end{tabular}

\vfill
{\footnotesize\color{finixgrey} Every repository path cited in this document exists in the FINIX repository.}
\end{titlepage}

\tableofcontents
\newpage

% ================================================================ 1
\section{Executive Summary}

FINIX is a digital banking platform built from twelve independently deployable services, each
owning its own logical database and communicating over HTTP and an event backbone. The platform
covers the full retail banking path: \textbf{identity} and customer profiles, \textbf{accounts}
with a reserve/commit balance model, \textbf{transaction orchestration} as a compensating saga, a
\textbf{cryptographic ledger} with double-entry postings and a tamper-evident hash chain,
\textbf{Vault} key custody with a 3-of-5 Shamir ceremony, \textbf{USSD} low-bandwidth banking,
\textbf{loan} processing, \textbf{compliance} case management and sanctions screening, an
\textbf{enclave runtime} for key reconstruction, \textbf{ML-assisted risk scoring},
a \textbf{payment hub} emitting ISO 20022 pacs.008, and \textbf{notifications}.

Phase 3 shifts the question from ``does it work?'' to ``does it keep working, and can you prove
it?'' FINIX answers that with five operational commitments, each backed by artifacts in the
repository: \textbf{deployment} of the exact images CI built for a given commit;
\textbf{automation} spanning build, publish, provision and release; \textbf{reliability} through
saga compensation, derived idempotency keys and an independently verifiable ledger;
\textbf{observability} through a metrics, logs and tracing stack with alert rules; and
\textbf{security engineering} across JWT identity, RFC 9449 DPoP verification, post-quantum
anchor signatures, and development-only gating of sensitive endpoints.

The platform runs live on a provisioned host behind a TLS-terminating reverse proxy:

\begin{center}
\begin{tabular}{@{}ll@{}}
\toprule
\textbf{Customer platform} & \url{https://roboti.qzz.io}\\
\textbf{Operations console} & \url{https://admin.roboti.qzz.io}\\
\bottomrule
\end{tabular}
\end{center}

% ================================================================ 2
\section{System Architecture}

FINIX is layered so that a failure or a redeploy in one tier does not require coordinated change
in another. Channels reach a single TLS edge; the edge reverse-proxies to services; services own
their data and publish events; the observability tier watches all of it.

\begin{center}
\fbox{\parbox{0.95\textwidth}{\vspace{4pt}\centering
{\small\textbf{\color{finixnavy}CHANNELS}}\\[3pt]
\fbox{\small Web PWA}\quad\fbox{\small Admin Console}\quad\fbox{\small Offline Vouchers}\quad\fbox{\small USSD}\quad\fbox{\small Lite ($<$50\,KB)}\\[6pt]
$\downarrow$\\[4pt]
{\small\textbf{\color{finixnavy}EDGE}}\\[3pt]
\fbox{\small Caddy --- automatic TLS}\quad\fbox{\small web nginx}\quad\fbox{\small admin nginx}\quad\fbox{\small Keycloak}\\[6pt]
$\downarrow$\\[4pt]
{\small\textbf{\color{finixnavy}SERVICES}}\\[3pt]
\fbox{\small Identity}\;\fbox{\small Account}\;\fbox{\small Ledger}\;\fbox{\small Orchestrator}\;\fbox{\small Vault}\;\fbox{\small USSD Gateway}\\[3pt]
\fbox{\small Loan}\;\fbox{\small Compliance}\;\fbox{\small Enclave Runtime}\;\fbox{\small Risk AI}\;\fbox{\small Payment Hub}\;\fbox{\small Notification}\\[6pt]
$\downarrow$\\[4pt]
{\small\textbf{\color{finixnavy}DATA / EVENT / SECURITY}}\\[3pt]
\fbox{\small PostgreSQL --- database per service}\quad\fbox{\small Redis}\quad\fbox{\small Redpanda (Kafka API)}\quad\fbox{\small Keycloak realm}\quad\fbox{\small OPA}\\[6pt]
$\downarrow$\\[4pt]
{\small\textbf{\color{finixnavy}OBSERVABILITY}}\\[3pt]
\fbox{\small Prometheus}\;\fbox{\small Alertmanager}\;\fbox{\small Grafana}\;\fbox{\small Loki}\;\fbox{\small Alloy}\;\fbox{\small Tempo}\\[3pt]
\fbox{\small node-exporter}\;\fbox{\small cAdvisor}\;\fbox{\small postgres-exporter}\;\fbox{\small redis-exporter}\;\fbox{\small blackbox-exporter}
\vspace{5pt}}}
\end{center}

\noindent\textbf{Architectural decisions are recorded, not implied.} The repository carries
Architecture Decision Records at \pth{docs/adr/}, including a logical database per service
(ADR-0003), post-quantum ML-DSA anchor signatures instead of BLS (ADR-0004), a transactional
outbox rather than change-data-capture (ADR-0005), and deploying published images rather than
server-side builds (ADR-0007).

\shot{Production architecture / service topology}
{Service topology as deployed, showing the edge, the running services and the data tier}
{\pth{docker compose ps} on the production host, or the Grafana service-overview dashboard}

% ================================================================ 3
\section{Phase 3 Rubric Evidence}

\subsection{A. Service Deployment \& Environment Consistency \hfill \rubric{15\%}}

\subsubsection{What FINIX does}
Every service ships as a container image built from a per-service \pth{Dockerfile}. The same
Compose model describes the stack in every environment, with a production overlay applied on top,
so a service behaves identically on a developer machine and on the production host. Images are
addressed by an immutable \pth{sha-<short>} tag derived from the commit, which makes the deployed
artifact traceable to a single revision. Service startup order is governed by health conditions
rather than timing assumptions.

\subsubsection{Repository evidence}
\begin{center}
\begin{tabularx}{\textwidth}{@{}lX@{}}
\toprule
\textbf{Artifact} & \textbf{Path}\\
\midrule
Base stack model & \pth{infra/compose/docker-compose.yml}\\
Production overlay & \pth{infra/compose/docker-compose.prod.yml}\\
Scale overlay & \pth{infra/compose/docker-compose.scale.yml}\\
Per-service images & \pth{services/*/Dockerfile}\\
Database bootstrap & \pth{infra/compose/init-databases.sh}\\
Local parity entrypoint & \pth{scripts/up.sh}, \pth{scripts/up-pull.sh}\\
\bottomrule
\end{tabularx}
\end{center}

\subsubsection{Live proof}
\shot{Running deployment}{\pth{docker compose ps} listing services and health state on the production host}{Terminal output on the server showing service names, status and health}

\subsection{B. Build \& Release Automation \hfill \rubric{20\%}}

\subsubsection{What FINIX does}
Three GitHub Actions workflows carry a commit from push to production. Continuous integration
compiles and tests the JVM modules and additionally runs the Go, Python and Node test suites and a
supply-chain scan. A publish workflow builds and pushes each service image to GHCR tagged with the
commit SHA. The deploy workflow then releases \emph{that exact image set} to the server. This is a
deliberate design decision recorded in ADR-0007: production runs the artifact CI built and tested
for the commit, rather than a rebuild performed on the server. A nightly chaos workflow exercises
failure behaviour on a schedule.

\subsubsection{Repository evidence}
\begin{center}
\begin{tabularx}{\textwidth}{@{}lX@{}}
\toprule
\textbf{Stage} & \textbf{Path}\\
\midrule
Continuous integration & \pth{.github/workflows/ci.yml}\\
Image publication & \pth{.github/workflows/publish-images.yml}\\
Release to production & \pth{.github/workflows/deploy.yml}\\
Scheduled chaos & \pth{.github/workflows/nightly-chaos.yml}\\
Server-side release & \pth{infra/deploy/deploy.sh}\\
Release rollback & \pth{infra/deploy/rollback.sh}\\
Decision record & \pth{docs/adr/0007-deploy-published-images-not-server-builds.md}\\
\bottomrule
\end{tabularx}
\end{center}

\subsubsection{Live proof}
\shot{GitHub Actions --- successful build/release pipeline}{Workflow summary showing green jobs and the commit SHA}{Actions run page for the released commit: CI, publish-images and deploy all green}

\subsection{C. Automated Infrastructure \& Configuration Management \hfill \rubric{15\%}}

\subsubsection{What FINIX does}
Host provisioning is expressed as Ansible roles rather than manual steps: base hardening and swap,
Docker installation, the Caddy reverse proxy with automatic TLS, the FINIX application role, and a
firewall role. Re-running the playbook is safe because the tasks are written to be idempotent.
Application credentials are generated on the host and mounted as Docker secrets rather than being
committed, with helper scripts for generation and rotation.

\subsubsection{Repository evidence}
\begin{center}
\begin{tabularx}{\textwidth}{@{}lX@{}}
\toprule
\textbf{Concern} & \textbf{Path}\\
\midrule
Playbook entrypoint & \pth{infra/ansible/site.yml}\\
Roles & \pth{infra/ansible/roles/} (base, docker, caddy, finix, firewall)\\
Inventories & \pth{infra/ansible/inventory/}\\
Secret generation & \pth{scripts/gen-secrets.sh}\\
Credential rotation & \pth{scripts/rotate-db-passwords.sh}\\
Decision record & \pth{docs/adr/0006-generated-docker-secrets-with-vault-as-runtime-source.md}\\
\bottomrule
\end{tabularx}
\end{center}

\subsubsection{Live proof}
\shot{Infrastructure automation}{Ansible playbook run against the production host, or the resulting host state}{Playbook recap showing changed/ok task counts, or \pth{systemctl}/\pth{docker} state after provisioning}

\subsection{D. Scalability, Availability \& Reliability \hfill \rubric{10\%}}

\subsubsection{What FINIX does}
Services are stateless at the request layer and hold state in PostgreSQL, so an instance can be
replaced without losing work. A scale overlay exists for running additional replicas of a service.
Availability is protected by container health checks and dependency-aware startup; reliability is
protected in the application layer by saga compensation, derived idempotency keys, downstream call
timeouts and resilience policies. Failure behaviour is exercised deliberately rather than assumed:
the repository contains a chaos scenario that kills the ledger mid-saga and records the result.

\subsubsection{Repository evidence}
\begin{center}
\begin{tabularx}{\textwidth}{@{}lX@{}}
\toprule
\textbf{Mechanism} & \textbf{Path}\\
\midrule
Replica overlay & \pth{infra/compose/docker-compose.scale.yml}\\
Chaos scenario \& results & \pth{tests/chaos/ledger-down.sh}, \pth{tests/chaos/RESULTS.md}\\
Contract tests & \pth{tests/contract/}\\
Load profile & \pth{tests/load/}\\
Saga orchestration & \pth{services/transaction-orchestrator/src/main/kotlin/}\\
Disaster-recovery runbook & \pth{docs/runbooks/dr-failover.md}\\
\bottomrule
\end{tabularx}
\end{center}

\subsubsection{Live proof}
\shot{Reliability / recovery evidence}{A dependency taken down mid-transaction and the platform's recorded response}{Chaos scenario output plus the saga's final state --- no stranded hold, no false COMPLETED}

\subsection{E. Operational Visibility \& System Health \hfill \rubric{15\%}}

\subsubsection{What FINIX does}
Every JVM service exposes Spring Boot Actuator health and a Prometheus scrape endpoint; the Go,
Python and Node services expose equivalents. Prometheus scrapes services and four exporters,
evaluates alert rules and routes to Alertmanager. Alloy ships logs to Loki; Tempo receives OTLP
traces. Grafana provisions its datasources from the repository so dashboards come up wired.

\subsubsection{Repository evidence}
\begin{center}
\begin{tabularx}{\textwidth}{@{}lX@{}}
\toprule
\textbf{Component} & \textbf{Path}\\
\midrule
Scrape configuration & \pth{infra/monitoring/prometheus/prometheus.yml}\\
Alert rules & \pth{infra/monitoring/prometheus/rules/finix.yml}\\
Alert routing & \pth{infra/monitoring/alertmanager/alertmanager.yml}\\
Log shipping & \pth{infra/monitoring/alloy/}, \pth{infra/monitoring/loki/}\\
Tracing & \pth{infra/monitoring/tempo/tempo.yml}\\
Endpoint probing & \pth{infra/monitoring/blackbox/blackbox.yml}\\
Dashboards \& datasources & \pth{infra/monitoring/grafana/}\\
\bottomrule
\end{tabularx}
\end{center}

\subsubsection{Live proof}
\shot{Grafana system overview}{Service availability, request health and resource behaviour across the stack}{Grafana dashboard with live panels and a visible time range}

\subsection{F. Security \& Sensitive Data \hfill \rubric{15\%}}

\subsubsection{What FINIX does}
Security is implemented as a shared kernel rather than repeated per service: a single
auto-configured filter chain gives every service JWT resource-server validation against Keycloak,
realm-role mapping, stateless sessions and an RFC 9449 DPoP verification filter. Credentials are
generated on the host and mounted as Docker secrets, never committed. Sensitive development
endpoints --- the ledger tamper injector and the demo data seeders --- are gated so they are not
registered on an ordinary production boot. Section~\ref{sec:security} states the current
enforcement posture precisely.

\subsubsection{Repository evidence}
\begin{center}
\begin{tabularx}{\textwidth}{@{}lX@{}}
\toprule
\textbf{Control} & \textbf{Path}\\
\midrule
Shared security configuration & \pth{services/shared-kernel/src/main/kotlin/org/finix/kernel/security/}\\
DPoP verification filter & \pth{services/shared-kernel/src/main/kotlin/org/finix/kernel/security/DPoPFilter.kt}\\
Identity provider assets & \pth{infra/keycloak/}, \pth{infra/keycloak/spi/}\\
Secret generation / rotation & \pth{scripts/gen-secrets.sh}, \pth{scripts/rotate-db-passwords.sh}\\
Vault custody domain & \pth{services/vault-service/src/main/kotlin/org/finix/vault/domain/}\\
Enclave assets & \pth{infra/enclave/}\\
Key ceremony runbook & \pth{docs/runbooks/key-ceremony.md}\\
\bottomrule
\end{tabularx}
\end{center}

\subsubsection{Live proof}
\shot{Security scanning result}{Supply-chain and secret scanning stage from the CI run}{CI job output for the scanning step on the released commit}

\subsection{G. Engineering Best Practices \hfill \rubric{5\%}}

\subsubsection{What FINIX does}
The codebase follows a hexagonal layout --- \pth{domain}, \pth{application}, \pth{adapter} --- with
architecture tests enforcing the boundaries, static analysis and a coverage gate wired into the
Gradle \pth{verify} task. Contribution process is documented, code ownership is declared, issue and
pull-request templates are provided, and a pre-commit hook is checked in. Design decisions are
captured as ADRs and operational procedures as runbooks.

\subsubsection{Repository evidence}
\begin{center}
\begin{tabularx}{\textwidth}{@{}lX@{}}
\toprule
\textbf{Practice} & \textbf{Path}\\
\midrule
Contribution guide & \pth{CONTRIBUTING.md}\\
Code ownership & \pth{.github/CODEOWNERS}\\
Review scaffolding & \pth{.github/pull\_request\_template.md}, \pth{.github/ISSUE\_TEMPLATE/}\\
Pre-commit hook & \pth{.githooks/pre-commit}\\
Decision records & \pth{docs/adr/}\\
Runbooks & \pth{docs/runbooks/}\\
Evidence map & \pth{docs/EVIDENCE.md}\\
User guide & \pth{docs/USER-GUIDE.md}\\
\bottomrule
\end{tabularx}
\end{center}

\subsection{H. Team Contributions \hfill \rubric{5\%}}

Contribution detail is recorded in Section~\ref{sec:team}, mapped to repository evidence so each
claim is checkable against commit history.

% ================================================================ 4
\section{Production Deployment}

Production work is performed \textbf{on the server}, from artifacts that continuous integration
already built and tested. No production image is produced on a developer machine.

\subsection{Release path}

\begin{center}
\begin{tabularx}{\textwidth}{@{}p{4mm}lX@{}}
\toprule
& \textbf{Stage} & \textbf{What happens}\\
\midrule
1 & Source revision & A commit on the release branch fixes the exact revision to be deployed.\\
2 & Build \& test & \pth{.github/workflows/ci.yml} compiles and tests the JVM modules and runs the Go, Python and Node suites plus supply-chain scanning.\\
3 & Publish & \pth{.github/workflows/publish-images.yml} builds each service image and pushes it to GHCR tagged \pth{sha-<short>} for that commit.\\
4 & Release & \pth{.github/workflows/deploy.yml} invokes \pth{infra/deploy/deploy.sh} on the host, which pulls the \pth{sha-<short>} image set and brings the stack up with \pth{--no-build}, recording the release.\\
5 & Health verification & Container health checks and service health endpoints are polled until the stack reports healthy.\\
6 & Smoke verification & A read-only smoke pass over the public edge confirms the customer and operations surfaces respond.\\
7 & Revision record & The deployed revision is recorded from the host for the evidence pack.\\
\bottomrule
\end{tabularx}
\end{center}

\subsection{Rollback}

Rollback is implemented as an ordinary deploy of an earlier recorded release:
\pth{infra/deploy/rollback.sh} reads the release history, resolves the target revision and
redeploys its published image set from GHCR --- a pull, not a rebuild. It accepts either the
previous release or an explicit revision, and can list the release history.

\subsection{Deployment evidence}

\shot{Exact deployed revision}{The commit SHA the running stack was built from}{\pth{git rev-parse HEAD} on the production host, alongside the matching image tag}

\shot{Running deployment / Compose services}{Service inventory and health as reported by the container runtime}{\pth{docker compose ps} on the production host}

\shot{Live customer dashboard}{The customer platform served over TLS at \url{https://roboti.qzz.io}}{Browser window showing the customer landing experience and a valid certificate}

\shot{Live operations dashboard}{The operations console served over TLS at \url{https://admin.roboti.qzz.io}}{Browser window showing the admin console}

\shot{CI workflow and release evidence}{The pipeline run that produced the deployed images}{Actions summary plus the GHCR package listing showing the \pth{sha-<short>} tag}

\noindent Final deployment evidence: [SCREENSHOT PLACEHOLDER --- see Section~\ref{sec:evidence}].

% ================================================================ 5
\section{Reliability \& Money Safety}

A banking platform is judged on what happens when something fails midway. FINIX treats a transfer
as a \textbf{saga}: a sequence of steps, each with a defined compensating action, whose progress is
persisted so that a restart resumes from fact rather than from memory.

\subsection{Transaction integrity}

\begin{center}
\begin{tabularx}{\textwidth}{@{}lX@{}}
\toprule
\textbf{Mechanism} & \textbf{Design}\\
\midrule
Saga orchestration & A transfer advances through persisted states; each downstream effect is recorded before the next is attempted.\\
Compensation & Every forward step has an inverse. A failure after a debit releases or refunds rather than leaving funds stranded.\\
Reserve / commit & Funds are reserved as a hold, then committed or released. A balance is never decremented speculatively.\\
Idempotency & Mutating downstream calls carry an idempotency key so a retried request is recognised as the same logical operation rather than a new one.\\
Fail-closed risk & A transfer proceeds only on an explicit risk decision; an indeterminate outcome does not become an implicit approval.\\
Double-entry ledger & Every posting balances debits against credits; an unbalanced journal entry is rejected.\\
Tamper-evident chain & Each journal entry carries the hash of its predecessor, so altering history breaks the chain at a determinable point.\\
Merkle anchoring & Entry windows are summarised into a Merkle root and signed, yielding inclusion proofs verifiable independently of the service.\\
Transactional outbox & Events are written in the same transaction as the state change, then published --- no lost or phantom events (ADR-0005).\\
Timeouts \& health checks & Downstream calls are bounded; containers declare health so an unhealthy instance is visible.\\
\bottomrule
\end{tabularx}
\end{center}

\subsection{Independent verification}

The ledger can be checked without trusting the service that produced it. \pth{scripts/verify-ledger.sh}
walks the chain through the public verification endpoint, recomputes what it expects and reports
the first break if one exists, then prints recent anchors. Inclusion proofs are Merkle paths that a
judge can fold independently.

\shot{Successful LKR 1 transfer}{A completed transfer between demo accounts, end to end}{Transfer confirmation in the customer UI plus the resulting saga state}

\shot{Ledger verification after transfer}{Chain verification reporting a valid ledger, and the inclusion proof for the transfer}{Output of \pth{scripts/verify-ledger.sh}, or the verification page showing the chain and proof}

\shot{Risk decision --- ALLOW / STEP-UP / BLOCK}{The three risk outcomes distinguished by the scoring service}{Three transfers producing an allow, a step-up challenge and a block, with the reason codes shown}

% ================================================================ 6
\section{Security Architecture}
\label{sec:security}

\subsection{Identity and session handling}

Keycloak is the identity provider, holding the \pth{finix} realm, its clients and its roles.
Services are configured as OAuth2 resource servers by the shared security kernel, which validates
JWTs against the issuer and maps realm roles onto Spring authorities. Browser clients use an
authorization-code flow with PKCE, and the web tier keeps token material in HttpOnly cookies rather
than exposing it to page scripts. A Keycloak authenticator SPI is included at \pth{infra/keycloak/spi/}
for login-risk enrichment.

\subsection{Sender-constrained tokens (DPoP)}

RFC 9449 DPoP verification capability is \textbf{implemented in the shared security layer and can be
enabled per service}. \pth{DPoPFilter} verifies proof signature, the \pth{htm} and \pth{htu}
bindings, the \pth{ath} access-token hash, proof freshness within a configured window, and
\pth{jti} uniqueness against a replay store; where an access token carries a \pth{cnf.jkt}
confirmation claim, the proof key thumbprint must match it. The filter is configuration-driven:
\pth{finix.dpop.enabled} controls whether it runs at all and \pth{finix.dpop.required} controls
whether a proof is mandatory for authenticated requests.

\noindent\textbf{Current posture, stated precisely:} the services ship with
\pth{finix.security.permit-all: true} and \pth{finix.dpop.required: false} in their
\pth{application.yml}, which is the demonstration configuration. Request-path authentication and
mandatory DPoP proofs are therefore \emph{available and configured but not currently enforced} in
the deployed environment; the enforcement switch is a configuration change, not new code. This
document does not claim otherwise.

\subsection{Offline value transfer}

Offline vouchers are signed on-device with ECDSA over a registered public key. Replay is prevented
by three independent controls: a monotonic per-device sequence number, a single-use nonce, and a
validity window. A device presenting a duplicate sequence or a reused nonce is rejected and
quarantined, so the same voucher cannot be spent twice when connectivity returns.

\subsection{Key custody --- demo-scale enclave boundary}

The master key is split with Shamir secret sharing over GF(256) into five custodian shards with a
3-of-5 reconstruction threshold, and Feldman verifiable secret sharing lets a shard be validated
against public commitments before use. Shards are sealed with a hybrid X25519 + ML-KEM-768 scheme.
Reconstruction is performed inside the enclave runtime, which returns the decrypted network
configuration rather than the master key itself, and every reconstruction appends to an egress log.
Ledger anchors are signed with post-quantum \textbf{ML-DSA} (ADR-0004).

\noindent This is a \textbf{demo-scale enclave boundary}. The custodians are five logical
identities rather than five physical devices, and the attestation document is a Nitro-format
document signed by a build-time key rather than output from Nitro hardware. The cryptography ---
Shamir, Feldman VSS, hybrid sealing and ML-DSA signing --- is genuinely implemented and is
exercised by the test suite.

\subsection{Development-only endpoints}

Two categories of endpoint exist for demonstration and must never be reachable in production:

\begin{itemize}
  \item \textbf{Ledger tamper injector} --- deliberately corrupts an entry so that chain
        verification can be shown detecting it. It is annotated \pth{@Profile("dev")} specifically,
        rather than including Spring's reserved \pth{default} profile, so a service booted with no
        profile selected never registers it. A dedicated test guards this.
  \item \textbf{Demo data seeders} --- create the persona accounts and profiles used in the
        walkthrough. These are separated into dedicated seeding beans under the same
        development-only gating, so production returns not-found rather than depending on an
        authorization rule to hold.
\end{itemize}

\subsection{Secret management}

No credential is committed. Secrets are generated on the host by \pth{scripts/gen-secrets.sh},
mounted into containers as Docker secrets, and rotatable via \pth{scripts/rotate-db-passwords.sh}
(ADR-0006). Supply-chain and secret scanning run as part of continuous integration.

\shot{Vault 3-of-5 ceremony}{A key ceremony reaching threshold and reconstructing inside the enclave}{Admin console ceremony view showing custodian approvals, threshold reached, and a non-empty egress log}

\shot{Security / secret scanning}{The scanning stage of the pipeline reporting on the released commit}{CI job output for supply-chain and secret scanning}

% ================================================================ 7
\section{Observability}

\subsection{What is instrumented}

Each JVM service exposes Actuator health and a Prometheus endpoint; the Go, Python and Node
services expose their own health and metrics surfaces. Prometheus scrapes those targets plus
node-exporter (host), cAdvisor (containers), postgres-exporter (database) and redis-exporter
(cache), with blackbox-exporter probing endpoints from the outside. Alert rules are evaluated in
Prometheus and routed through Alertmanager. Alloy collects container logs into Loki, and Tempo
receives OTLP traces from the services. Grafana provisions its datasources from the repository, so
metrics, logs and traces are queryable from one place.

\subsection{What a judge can inspect}

\begin{center}
\begin{tabularx}{\textwidth}{@{}lX@{}}
\toprule
\textbf{Question} & \textbf{Where it is answered}\\
\midrule
Is every service up? & Prometheus targets page; Grafana availability panels; container health state\\
Is request handling healthy? & Request rate, error rate and latency panels per service\\
How is the host behaving? & node-exporter and cAdvisor panels --- CPU, memory, container restarts\\
What happened, and when? & Loki log search across services, correlated by trace and correlation id\\
Did it recover? & Availability panel across the incident window, alongside the chaos scenario record\\
\bottomrule
\end{tabularx}
\end{center}

\shot{Grafana system overview}{Cross-service availability and request health}{Grafana dashboard with populated panels and a visible time range}

\shot{Prometheus targets / health}{Scrape targets and their up state}{Prometheus targets page showing services and exporters as up}

\shot{Logs / operational visibility}{Aggregated logs for a real request path}{Loki query for a transfer, showing correlated entries across services}

% ================================================================ 8
\section{Resilient Banking Channels}

Financial inclusion is a bandwidth and device problem as much as a software problem. FINIX
therefore serves four distinct access conditions from the same core.

\begin{center}
\begin{tabularx}{\textwidth}{@{}llX@{}}
\toprule
\textbf{Channel} & \textbf{Condition served} & \textbf{Design}\\
\midrule
Web PWA & Modern device, normal connectivity & Full customer experience: transfers, verification, history\\
Offline vouchers & Intermittent or absent connectivity & Value transfer signed on-device and reconciled when connectivity returns, protected against replay\\
USSD & Feature phone, no data plan & Menu-driven balance and transfer flows over a session protocol, with a simulator for demonstration\\
Lite interface & Severely constrained bandwidth & A budgeted static page held under 50\,KB with zero client JavaScript, enforced by \pth{scripts/check-lite-budget.sh}\\
\bottomrule
\end{tabularx}
\end{center}

\noindent The lite budget is not aspirational --- it is a build-time check that fails if the page
exceeds its size ceiling or introduces a script tag.

\shot{Customer dashboard}{The main customer experience on the live platform}{Browser at \url{https://roboti.qzz.io} showing the customer dashboard}

\shot{Offline mode}{Voucher creation or reconciliation while disconnected}{Offline flow with the network disabled, then the reconciliation result}

\shot{USSD}{A USSD session producing a balance or transfer}{USSD simulator showing the menu and the session response}

\shot{Lite page}{The low-bandwidth interface and its transferred size}{Lite page with browser devtools showing transferred bytes under the budget}

% ================================================================ 9
\section{Live Demonstration Plan}

The walkthrough is designed to fit inside seven minutes and to show running behaviour rather than
slides.

\begin{center}
\begin{longtable}{@{}p{11mm}p{31mm}p{40mm}p{40mm}p{10mm}@{}}
\toprule
\textbf{Time} & \textbf{Action} & \textbf{What the judge sees} & \textbf{What it proves} & \textbf{Rubric}\\
\midrule
\endfirsthead
\toprule
\textbf{Time} & \textbf{Action} & \textbf{What the judge sees} & \textbf{What it proves} & \textbf{Rubric}\\
\midrule
\endhead
0:00 & Production dashboard & Live customer and admin URLs over TLS & The platform is deployed and publicly reachable & A\\
0:30 & CI / deployment evidence & Green pipeline, image tag, deployed revision on the host & The running system is the artifact CI built for that commit & B\\
1:15 & Live transfer & LKR 1 moved between demo accounts, confirmed & The money path works end to end & D\\
2:15 & Risk decision & Allow, step-up and block outcomes with reason codes & Risk is evaluated and fails closed & D, F\\
3:00 & Ledger verification & Chain verified, inclusion proof folded independently & Financial history is tamper-evident and independently checkable & D\\
3:45 & Reliability / recovery & A dependency killed mid-saga; compensation observed & Failure does not strand or duplicate money & D\\
4:30 & Grafana & Availability, request health, resource panels, logs & The platform is observable in operation & E\\
5:20 & Vault ceremony & 3-of-5 approvals, reconstruction, egress log & Key custody is threshold-based and auditable & F\\
6:00 & Offline / USSD / Lite & Three constrained channels exercised & Access is resilient across devices and bandwidth & D, G\\
6:40 & Architecture \& close & Topology, ADRs, evidence pack & The design is deliberate and documented & G, H\\
\bottomrule
\end{longtable}
\end{center}

% ================================================================ 10
\section{Screenshot Evidence}
\label{sec:evidence}

The following numbered placeholders are the evidence pack. Each compiles without an image file
present and is replaced by the corresponding capture before submission.

\evidence{1}{Customer production dashboard}{Browser at \url{https://roboti.qzz.io} with a valid TLS certificate visible}
\evidence{2}{Admin operations dashboard}{Browser at \url{https://admin.roboti.qzz.io} showing the operations console}
\evidence{3}{GitHub Actions successful pipeline}{Workflow summary with green CI, publish and deploy jobs and the commit SHA}
\evidence{4}{Running deployment / Compose services}{\pth{docker compose ps} on the production host showing services and health}
\evidence{5}{Exact deployed revision}{\pth{git rev-parse HEAD} on the host next to the deployed image tag}
\evidence{6}{Live transfer}{A completed LKR 1 transfer with its confirmation and resulting state}
\evidence{7}{Ledger verification}{Chain verification result and a Merkle inclusion proof for the transfer}
\evidence{8}{Grafana}{System overview dashboard with populated panels and a visible time range}
\evidence{9}{Vault ceremony}{3-of-5 custodian approvals, reconstruction outcome and egress log}
\evidence{10}{Offline / USSD / Lite}{The three constrained channels, captured together or as a strip}
\evidence{11}{Security / scanning result}{Supply-chain and secret scanning output from the released commit}
\evidence{12}{Rollback / recovery evidence}{Release history and a redeploy of an earlier revision, if exercised}

% ================================================================ 11
\section{Team Contributions}
\label{sec:team}

\begin{center}
\begin{tabularx}{\textwidth}{@{}p{26mm}p{30mm}XX@{}}
\toprule
\textbf{Contributor} & \textbf{Area} & \textbf{Key contribution} & \textbf{Repository evidence}\\
\midrule
{[NAME]} & [AREA] & [CONTRIBUTION] & [PATH / PR / COMMIT]\\
{[NAME]} & [AREA] & [CONTRIBUTION] & [PATH / PR / COMMIT]\\
{[NAME]} & [AREA] & [CONTRIBUTION] & [PATH / PR / COMMIT]\\
{[NAME]} & [AREA] & [CONTRIBUTION] & [PATH / PR / COMMIT]\\
{[NAME]} & [AREA] & [CONTRIBUTION] & [PATH / PR / COMMIT]\\
\bottomrule
\end{tabularx}
\end{center}

\noindent Contribution claims are checkable against commit history and the code ownership map at
\pth{.github/CODEOWNERS}.

% ================================================================ 12
\section{Final Verification Checklist}

\begin{center}
\begin{tabularx}{\textwidth}{@{}p{7mm}X@{}}
\toprule
& \textbf{Item}\\
\midrule
\chk & Repository publicly accessible\\
\chk & Live customer URL reachable over TLS --- \url{https://roboti.qzz.io}\\
\chk & Live operations URL reachable over TLS --- \url{https://admin.roboti.qzz.io}\\
\chk & Final intended commit deployed to production\\
\chk & Continuous integration green on the deployed commit\\
\chk & All critical services reporting healthy\\
\chk & Read-only smoke verification passed against the public edge\\
\chk & Live transfer verified end to end\\
\chk & Ledger chain verified and an inclusion proof checked\\
\chk & Grafana evidence captured\\
\chk & Vault ceremony evidence captured\\
\chk & Development-only endpoints confirmed unavailable in production\\
\chk & No credentials present in the repository\\
\chk & Deployment screenshots captured and inserted\\
\chk & Team contribution table completed\\
\chk & Submission links checked\\
\bottomrule
\end{tabularx}
\end{center}

% ================================================================ 13
\section{Conclusion}

FINIX demonstrates that financial correctness and operational discipline are the same engineering
problem approached from two directions. The platform combines \textbf{financial correctness} ---
double-entry postings, a tamper-evident hash chain, Merkle inclusion proofs and saga compensation
that refuses to strand or duplicate money; \textbf{independent services} --- twelve deployable
units, each owning its data, composed rather than entangled; \textbf{automation} --- a commit that
travels to production as the exact artifact continuous integration tested, on infrastructure
described as code; \textbf{observability} --- metrics, logs, traces and alert rules that make the
running system answerable; \textbf{resilient access} --- web, offline, USSD and a sub-50\,KB lite
channel so that connectivity is not a precondition for banking; and \textbf{security engineering}
--- threshold key custody, post-quantum anchor signatures, sender-constrained token verification,
and demonstration endpoints deliberately gated out of production boots.

\vspace{4mm}
\begin{center}
{\Large\itshape\color{finixnavy} Deploy It. Automate It. Keep It Running.}
\end{center}
\vspace{4mm}

\begin{center}
\fbox{\parbox{0.93\textwidth}{\vspace{3pt}
\textbf{Evaluation note:} FINIX should be assessed against the repository artifacts, live system
behaviour, and evidence shown in this document rather than promotional wording alone.
\vspace{3pt}}}
\end{center}

\end{document}
